% MT049 / explanation — canonical editable content.
题目同时提出“预测要及时”和“新增数据后还要更新”。第一件事关乎模型执行成本，第二件事关乎训练机制，应该分别回答。


\textbf{(a) 控制每轮成本与总轮数。} 减少弱学习器数量、限制它们的复杂度，可以减少训练与预测工作。学习率控制新学习器加入组合的力度，适当增大可能用较少轮数取得效果，但步子过大也会使训练不稳定或过度追逐噪声。因此选择要落到验证数据上的精度、时延与模型大小，不能只把学习率无限调大。


\textbf{(b) 常规批量AdaBoost不自动适合持续在线更新。} 第\(t\)轮训练使用前面结果得到的样本权重。新数据到来后，怎样与旧权重结合、是否重训已有学习器，都需要额外设计；“每次追加一个学习器”并不是无成本的更新方案，追加过多还会使预测越来越慢。


AdaBoost的指数损失是\(e^{-yf(x)}\)。某点严重错分时，\(yf(x)\)很负，损失与梯度会很大，模型会努力纠正它；若该点其实是传感噪声或错误标签，这种机制就可能浪费拟合能力。损失图的解释见\StudyRef{Answers}{mt065}{MT065}。


因此可建议限制模型大小、定期用经过检查的数据重训，或选择专门支持流式更新的方法。这里评价的是题述批量方案；训练好的AdaBoost仍可以并行预测，也不能据此断言所有在线AdaBoost变体都不可行。


\textbf{出处：} 2023A期中 Q11，原卷第8页（10分）


\textbf{答案依据：} 2023A期中 Q11 随卷／配套参考答案；条件与解释经复核。
